% Generated draft: requires mathematical and source review.
\documentclass[11pt]{article}
\usepackage[margin=25mm]{geometry}
\usepackage{fontspec}
\setmainfont{Latin Modern Roman}
\usepackage{amsmath,amssymb,mathtools,unicode-math}
\setmathfont{Latin Modern Math}
\usepackage{enumitem,needspace,xurl,hyperref}
\setlength{\parindent}{0pt}
\setlength{\parskip}{5pt}
\setlength{\emergencystretch}{3em}
\newcommand{\N}{\mathbb{N}}
\newcommand{\Z}{\mathbb{Z}}
\newcommand{\Q}{\mathbb{Q}}
\newcommand{\R}{\mathbb{R}}
\newcommand{\C}{\mathbb{C}}
\newcommand{\F}{\mathbb{F}}
\newcommand{\A}{\mathbb{A}}
\newcommand{\PP}{\mathbb{P}}
\newcommand{\E}{\mathbb{E}}
\newcommand{\eps}{\varepsilon}
\newcommand{\dd}{\,\mathrm{d}}
\begin{document}
\section*{UnsolvedMath ID: 30006238}
\textbf{Source problem code:} catalogue entry
\subsection{Definitions and mathematical statement}
Let \(\mathscr L_{\mathrm{src}}\) denote the class of objects called localities in the cited report. For \(\mathcal L\in\mathscr L_{\mathrm{src}}\), write \(\operatorname{Simple}_{\mathrm{src}}(\mathcal L)\), \(\operatorname{Proper}_{\mathrm{src}}(\mathcal L)\), and \(\operatorname{LF}_{\mathrm{src}}(\mathcal L)\) for the report's predicates of being simple, proper, and locally finite, respectively.

The question is whether
\[
\exists\,\mathcal L\in\mathscr L_{\mathrm{src}}\quad
\operatorname{Simple}_{\mathrm{src}}(\mathcal L)
\;\land\;
\operatorname{Proper}_{\mathrm{src}}(\mathcal L)
\;\land\;
\neg\operatorname{LF}_{\mathrm{src}}(\mathcal L).
\]
Thus, does there exist a simple proper locality that is not locally finite?

These predicates retain their source-specific technical meanings. The supplied record does not define a locality, simplicity, properness, or local finiteness in this setting, or specify the ambient restrictions and prime conventions. Consequently, this is a source-relative formulation; a fully self-contained mathematical statement requires those definitions and hypotheses from the report. In particular, no finiteness assumption on the locality, or restriction to compact localities, is added.
\subsection{Short English statement}
Does a simple proper locality exist that is not locally finite, using the cited report's definitions? The supplied record labels this question open, but its technical definitions and current literature status require independent review.
\subsection{Sources}
\begin{itemize}
\item Supplied local UnsolvedMath registry record: permanent ID \(30006238\), title \emph{Non-Locally-Finite Simple Proper Localities}, source problem code \texttt{catalogue entry}, in the Oberwolfach Reports Open Problems collection. The record gives proposed year \(2025\) and supplies no named proposer.
\item \emph{Finite Groups, Fusion Systems and Applications} (2025), Oberwolfach Reports, pp.~689--752, DOI \texttt{10.4171/owr/2025/16}. Supplied source URL: \url{https://doi.org/10.4171/owr/2025/16}. No more precise location within the report is supplied.
\item Supplied literature reference, labelled ``Gonzales locality results and the stated open question'': \url{https://doi.org/10.4171/owr/2025/16}. This label does not by itself identify the proposer of the question.
\item Literature assessment dated exactly \texttt{2026-08-22}: the supplied curation labels the question open and states that the report establishes local finiteness for compact localities, without establishing it for the general case asked about here. This is supplied dated curation, not an independently verified literature assessment or verification of the cited mathematical claims.
\item The registry labels the statement \texttt{corrected\_verified}, with the qualification that only clear formula delimiters, notation, grammar, and typography were repaired and that mathematical content and scope were preserved. That supplied editorial label does not constitute mathematical review of this draft.
\end{itemize}
\end{document}
